% GENERATED FILE. Edit canonical lessons, then run npm run generate:books.
% canonical-source-hash: fnv1a64:51e46fa502a3b24a
% canonical-lessons: GU-C398-jhoku, GU-C398-ujagro, GU-C398-ughvu, GU-C398-jhatko, GU-C398-khakhadvu

\chapter{Dozing on the Train}
\label{ch:gu-dozing-on-the-train}
\hlchaptermodality{398}

\begin{chapteropening}
\textbf{By the end of this chapter:} \emph{I can say how I slept on a night journey.}

Complete the last lesson of chapter 398: I can say how I slept on a night journey.
\end{chapteropening}

\section[\texorpdfstring{\gu{ઝોકું} (jhokũ)}{ઝોકું (jhokũ)}]{\gu{ઝોકું} (jhokũ) — a nap; a nod of the head in sleep}

\label{lesson:GU-C398-jhoku}



\begin{quote}
Before the new one: say the Gujarati for to shove, then the Gujarati for to fill up.
\end{quote}

\subsection*{You'll want to know: \gu{ઝોકું}}
\textbf{\gu{ઝોકું}} — \emph{jhokũ} — \textquotedblleft{}a nap; a nod of the head in sleep\textquotedblright{}.

\begin{quote}
\gu{બપોરે} \gu{દાદાને} \gu{ટ્રેનમાં} \gu{ઝોકું} \gu{આવ્યું}.
\end{quote}

\begin{quote}
\emph{bapore dādāne ṭrenmā̃ jhokũ āvyũ.}
\end{quote}

\begin{quote}
In the afternoon grandfather nodded off on the train.
\end{quote}

A neuter noun.

\begin{cousinweb}
\textbf{\gu{ઝોકું}} (\emph{jhokũ}) seems to belong with the Gujarati verb \emph{jhukvũ}, \textquotedblleft{}to bend\textquotedblright{}, which CDIAL lists under \emph{*jhukkati}, \textquotedblleft{}stoops\textquotedblright{}, a form it reconstructs (CDIAL 5399): the head bends forward as you doze. CDIAL gives no form for \gu{ઝોકું} itself, so the link is hedged.
\end{cousinweb}

\begin{grammarlens}[title={What you've built}]
One more word for getting around.
\end{grammarlens}

\subsection*{Your turn}
\begin{itemize}
\raggedright
  \item \emph{Say it:} \emph{jhokũ}
  \item \emph{Say it:} \emph{jhokũ}, once more
  \item \emph{Say it:} \emph{bapore dādāne ṭrenmā̃ jhokũ āvyũ.} — In the afternoon grandfather nodded off on the train.
  \item \emph{Recall:} say the Gujarati for to shove, then the Gujarati for to fill up, then say \emph{jhokũ} again
\end{itemize}

\begin{tcolorbox}[breakable,colback=teal!4,colframe=teal!35!black,title={Before you move on}]
What does \gu{ઝોકું} mean? (\textquotedblleft{}A nap; a nod of the head in sleep\textquotedblright{}.) Say it once more.
\end{tcolorbox}
\section[\texorpdfstring{\gu{ઉજાગરો} (ujāgro)}{ઉજાગરો (ujāgro)}]{\gu{ઉજાગરો} (ujāgro) — a sleepless night}

\label{lesson:GU-C398-ujagro}



\begin{quote}
Before the new one: say the Gujarati for to fill up, then the Gujarati for a nap.
\end{quote}

\subsection*{You'll want to know: \gu{ઉજાગરો}}
\textbf{\gu{ઉજાગરો}} — \emph{ujāgro} — \textquotedblleft{}a sleepless night\textquotedblright{}.

\begin{quote}
\gu{આખી} \gu{રાત} \gu{ટ્રેનમાં} \gu{મને} \gu{ઉજાગરો} \gu{થયો}.
\end{quote}

\begin{quote}
\emph{ākhī rāt ṭrenmā̃ mane ujāgro thayo.}
\end{quote}

\begin{quote}
I had a sleepless night on the train.
\end{quote}

A masculine noun.

\begin{cousinweb}
\textbf{\gu{ઉજાગરો}} (\emph{ujāgro}) goes back to \emph{*ujjāgrat}, \textquotedblleft{}awake\textquotedblright{}, a form CDIAL reconstructs; CDIAL lists the Gujarati word, waking through the night (CDIAL 1665).
\end{cousinweb}

\begin{grammarlens}[title={What you've built}]
One more word for getting around.
\end{grammarlens}

\subsection*{Your turn}
\begin{itemize}
\raggedright
  \item \emph{Say it:} \emph{ujāgro}
  \item \emph{Say it:} \emph{ujāgro}, once more
  \item \emph{Say it:} \emph{ākhī rāt ṭrenmā̃ mane ujāgro thayo.} — I had a sleepless night on the train.
  \item \emph{Recall:} say the Gujarati for to fill up, then the Gujarati for a nap, then say \emph{ujāgro} again
\end{itemize}

\begin{tcolorbox}[breakable,colback=teal!4,colframe=teal!35!black,title={Before you move on}]
What does \gu{ઉજાગરો} mean? (\textquotedblleft{}A sleepless night\textquotedblright{}.) Say it once more.
\end{tcolorbox}
\section[\texorpdfstring{\gu{ઊંઘવું} (ū̃ghvũ)}{ઊંઘવું (ū̃ghvũ)}]{\gu{ઊંઘવું} (ū̃ghvũ) — to be asleep; to sleep soundly}

\label{lesson:GU-C398-ughvu}



\begin{quote}
Before the new one: say the Gujarati for a nap, then the Gujarati for a sleepless night.
\end{quote}

\subsection*{You'll want to know: \gu{ઊંઘવું}}
\textbf{\gu{ઊંઘવું}} — \emph{ū̃ghvũ} — \textquotedblleft{}to be asleep; to sleep soundly\textquotedblright{}.

\begin{quote}
\gu{ટ્રેન} \gu{ચાલતી} \gu{હતી} \gu{ત્યારે} \gu{બાળકો} \gu{ઊંઘતાં} \gu{હતાં}.
\end{quote}

\begin{quote}
\emph{ṭren cāltī hatī tyāre bāḷko ū̃ghtā̃ hatā̃.}
\end{quote}

\begin{quote}
While the train was moving the children were asleep.
\end{quote}

Past (he …): \gu{તે} \gu{આખી} \gu{રાત} \gu{ઊંઘ્યો}, \emph{te ākhī rāt ū̃ghyo}.

\begin{cousinweb}
\textbf{\gu{ઊંઘવું}} (\emph{ū̃ghvũ}) goes back to \emph{*uṅghati}, \textquotedblleft{}sleeps\textquotedblright{}, a form CDIAL reconstructs; CDIAL lists the Gujarati verb (CDIAL 1632). \gu{ઊંઘ}, sleep, is the noun beside it.
\end{cousinweb}

\begin{grammarlens}[title={What you've built}]
One more word for getting around.
\end{grammarlens}

\subsection*{Your turn}
\begin{itemize}
\raggedright
  \item \emph{Say it:} \emph{ū̃ghvũ}
  \item \emph{Say it:} \emph{ū̃ghvũ}, once more
  \item \emph{Say it:} \emph{ṭren cāltī hatī tyāre bāḷko ū̃ghtā̃ hatā̃.} — While the train was moving the children were asleep.
  \item \emph{Recall:} say the Gujarati for a nap, then the Gujarati for a sleepless night, then say \emph{ū̃ghvũ} again
\end{itemize}

\begin{tcolorbox}[breakable,colback=teal!4,colframe=teal!35!black,title={Before you move on}]
What does \gu{ઊંઘવું} mean? (\textquotedblleft{}To be asleep; to sleep soundly\textquotedblright{}.) Say it once more.
\end{tcolorbox}
\section[\texorpdfstring{\gu{ઝટકો} (jhaṭko)}{ઝટકો (jhaṭko)}]{\gu{ઝટકો} (jhaṭko) — a jolt; a jerk}

\label{lesson:GU-C398-jhatko}



\begin{quote}
Before the new one: say the Gujarati for a sleepless night, then the Gujarati for to be asleep.
\end{quote}

\subsection*{You'll want to know: \gu{ઝટકો}}
\textbf{\gu{ઝટકો}} — \emph{jhaṭko} — \textquotedblleft{}a jolt; a jerk\textquotedblright{}.

\begin{quote}
\gu{ટ્રેન} \gu{અચાનક} \gu{અટકી} \gu{અને} \gu{ડબ્બામાં} \gu{ઝટકો} \gu{લાગ્યો}.
\end{quote}

\begin{quote}
\emph{ṭren acānak aṭkī ane ḍabbāmā̃ jhaṭko lāgyo.}
\end{quote}

\begin{quote}
The train stopped suddenly and the carriage gave a jolt.
\end{quote}

A masculine noun.

\begin{cousinweb}
\textbf{\gu{ઝટકો}} (\emph{jhaṭko}) seems to be built on \gu{ઝટ}, \textquotedblleft{}quickly, at once\textquotedblright{}, which CDIAL lists for Gujarati under \emph{*jhaṭ}, a form it reconstructs (CDIAL 5327): a sudden quick shake.
\end{cousinweb}

\begin{grammarlens}[title={What you've built}]
One more word for getting around.
\end{grammarlens}

\subsection*{Your turn}
\begin{itemize}
\raggedright
  \item \emph{Say it:} \emph{jhaṭko}
  \item \emph{Say it:} \emph{jhaṭko}, once more
  \item \emph{Say it:} \emph{ṭren acānak aṭkī ane ḍabbāmā̃ jhaṭko lāgyo.} — The train stopped suddenly and the carriage gave a jolt.
  \item \emph{Recall:} say the Gujarati for a sleepless night, then the Gujarati for to be asleep, then say \emph{jhaṭko} again
\end{itemize}

\begin{tcolorbox}[breakable,colback=teal!4,colframe=teal!35!black,title={Before you move on}]
What does \gu{ઝટકો} mean? (\textquotedblleft{}A jolt; a jerk\textquotedblright{}.) Say it once more.
\end{tcolorbox}
\section[\texorpdfstring{\gu{ખખડવું} (khakhaḍvũ)}{ખખડવું (khakhaḍvũ)}]{\gu{ખખડવું} (khakhaḍvũ) — to rattle}

\label{lesson:GU-C398-khakhadvu}



\begin{quote}
Before the new one: say the Gujarati for to be asleep, then the Gujarati for a jolt.
\end{quote}

\subsection*{You'll want to know: \gu{ખખડવું}}
\textbf{\gu{ખખડવું}} — \emph{khakhaḍvũ} — \textquotedblleft{}to rattle\textquotedblright{}.

\begin{quote}
\gu{જૂની} \gu{બારી} \gu{આખો} \gu{દિવસ} \gu{ખખડતી} \gu{હતી}.
\end{quote}

\begin{quote}
\emph{jūnī bārī ākho divas khakhaḍtī hatī.}
\end{quote}

\begin{quote}
The old window rattled all day.
\end{quote}

Past (it …): \gu{તે} \gu{બારી} \gu{ખખડી}, \emph{te bārī khakhḍī}.

\begin{cousinweb}
\textbf{\gu{ખખડવું}} (\emph{khakhaḍvũ}) sounds like the noise it names. CDIAL lists a Gujarati verb of the same kind, \emph{khaṛkhaṛvũ}, \textquotedblleft{}to rattle\textquotedblright{}, under Sanskrit \emph{khaṭakhaṭāyatē}, \textquotedblleft{}crackles\textquotedblright{} (CDIAL 3771), beside Hindi \emph{khaṛkhaṛānā} and Marathi \emph{khaḍkhaḍṇẽ}; it gives no form for \gu{ખખડવું} itself, so the link is hedged.
\end{cousinweb}

\begin{grammarlens}[title={What you've built}]
One more word for getting around.
\end{grammarlens}

\subsection*{Your turn}
\begin{itemize}
\raggedright
  \item \emph{Say it:} \emph{khakhaḍvũ}
  \item \emph{Say it:} \emph{khakhaḍvũ}, once more
  \item \emph{Say it:} \emph{jūnī bārī ākho divas khakhaḍtī hatī.} — The old window rattled all day.
  \item \emph{Recall:} say the Gujarati for to be asleep, then the Gujarati for a jolt, then say \emph{khakhaḍvũ} again
\end{itemize}

\begin{tcolorbox}[breakable,colback=teal!4,colframe=teal!35!black,title={Before you move on}]
What does \gu{ખખડવું} mean? (\textquotedblleft{}To rattle\textquotedblright{}.) Say it once more.
\end{tcolorbox}
